\section{问题重述}

某物流分拣中心已知未来30天的逐小时进货量，每日安排5个连续8小时班次，每名工人同一天只参加一个班次。常规处理能力为每人每小时25件，且当日货物必须当日处理完。任务是：首先设计每日班次及人数，使每日用工最少；其次考虑0--12时进货必须在16时前处理完，且每名工人每班有1小时只能处理10件；最后在第二问基础上按月招工，使每名工人工作23天、连续工作不超过7天，并最小化招工人数。

问题的关键并非简单地用“日到货量/日处理量”估算人数，而是处理三个耦合关系：到货发生在不同小时，早期空闲产能不能服务尚未到达的货物；连续8小时班次产生跨小时容量耦合；月度固定工日与连续工作限制又把30个日模型连接为一个整体。

\section{数据分析与问题分析}

附件包含30天、每天下24个小时共720条数据。全月进货量为2155604件，小时均值为2993.89件，单小时最大值7371件；日到货总量在59156--107377件之间。图\ref{fig:heatmap}显示进货高峰随日期发生迁移，前半月多在深夜出现，后半月分布更平缓，因此固定人数、固定起点的排班方案会造成明显的峰谷错配。

\begin{figure}[htbp]
  \centering
  \includegraphics[width=.88\textwidth]{figures/arrival_heatmap.pdf}
  \caption{30天逐小时进货量热力图}
  \label{fig:heatmap}
\end{figure}

本文采用逐日滚动求解：问题一和问题二分别建立日混合整数规划，得到每天5个班次的起点与人数；问题三再将问题二的每日人数作为覆盖需求，建立工人--日期0--1排班模型，并把当日到岗人员分配回相应班次。该分层方式既保持日内货物流精确，又避免一次性模型规模过大。

\section{模型假设}

\begin{enumerate}[label=\arabic*)]
  \item 班次在整数小时开始，起点属于$S=\{0,1,\ldots,16\}$，持续8小时且不跨自然日；
  \item 同类工人处理效率相同，交接、设备故障和行走时间已包含在给定小时处理能力中；
  \item 货物到达后才能处理，未处理货物可在当日后续小时形成库存，但不允许跨日；
  \item 问题二中每名工人的低效率小时可在其8小时班次内错峰安排；
  \item 问题三允许实际到岗人数高于当日最低需求，多出的人员作为机动冗余并参加一个班次。
\end{enumerate}

\section{符号说明}

\begin{table}[htbp]
\centering
\caption{主要符号及含义}
\label{tab:symbols}
\begin{tabular}{cll}
\toprule
符号 & 含义 & 单位\\
\midrule
$a_{dh}$ & 第$d$天第$h$小时进货量 & 件\\
$x_{ds}$ & 第$d$天是否启用起点为$s$的班次 & 0--1\\
$y_{ds}$ & 第$d$天起点为$s$的班次人数 & 人\\
$p_{dh}$ & 第$d$天第$h$小时实际处理量 & 件\\
$b_{dh}$ & 第$d$天第$h$小时结束后的待处理量 & 件\\
$z_{dsk}$ & 班次$s$第$k$个小时处于低效率状态的人数 & 人\\
$w_{id}$ & 工人$i$在第$d$天是否工作 & 0--1\\
$r_d$ & 问题二给出的第$d$天最低人数 & 人\\
\bottomrule
\end{tabular}
\end{table}
